\documentclass[12pt,a4paper]{article}
\usepackage[utf8]{inputenc}
\usepackage[spanish,es-tabla,es-noquoting]{babel}
\usepackage{geometry}
\usepackage{graphicx}
\usepackage{multirow}
\usepackage{tabularx}
\usepackage{mathptmx}
\usepackage{fancyhdr}
\usepackage{array}
\usepackage[table]{xcolor}
\usepackage{titlesec}
\usepackage{setspace}
\usepackage{microtype}
\usepackage{lastpage}
\usepackage{booktabs}
\usepackage{enumitem}
\usepackage{amsmath}
\usepackage{amsfonts}
\usepackage{amssymb}
\usepackage{placeins}
\usepackage{float} % Posición rígida [H]
\usepackage{tikz}
\usepackage{pgfplots}
\pgfplotsset{compat=1.17}
\usetikzlibrary{arrows.meta,positioning,shapes.geometric,calc}

\graphicspath{{images/}{../images/}{./}}

\geometry{
    top=4.7cm,
    bottom=2.5cm,
    left=2.5cm,
    right=2.5cm,
    headheight=3.9cm,
    headsep=0.4cm
}

\setstretch{1.15}
\setlength{\parindent}{0pt}
\setlength{\parskip}{0.45em}
\setlength{\tabcolsep}{4pt}

\titleformat{\section}
  {\normalfont\fontsize{12}{14.5}\bfseries}{\thesection.}{0.6em}{\uppercase}
\titlespacing*{\section}
  {0pt}{1.8ex plus 0.4ex minus 0.2ex}{0.6ex plus 0.2ex}

\titleformat{\subsection}
  {\normalfont\fontsize{11.5}{13.5}\bfseries}{\thesubsection.}{0.5em}{}
\titlespacing*{\subsection}
  {0pt}{1.4ex plus 0.3ex minus 0.1ex}{0.4ex plus 0.1ex}

\titleformat{\subsubsection}
  {\normalfont\fontsize{10.5}{12.5}\bfseries\itshape}{\thesubsubsection.}{0.4em}{}
\titlespacing*{\subsubsection}
  {0pt}{1.1ex plus 0.2ex minus 0.1ex}{0.3ex plus 0.1ex}

\newcolumntype{Y}{>{\centering\arraybackslash}X}
\newcolumntype{L}{>{\raggedright\arraybackslash}X}
\newcolumntype{C}{>{\centering\arraybackslash}X}

\definecolor{headerred}{RGB}{180, 0, 0}
\definecolor{darkred}{RGB}{140, 0, 0}
\definecolor{lightgraybg}{RGB}{248, 248, 248}
\definecolor{tableheader}{RGB}{232, 238, 246}
\definecolor{navyheader}{RGB}{0, 32, 96}
\definecolor{lineblue}{RGB}{0, 102, 204}
\definecolor{areablue}{RGB}{220, 236, 252}
\definecolor{lcsred}{RGB}{217, 56, 41}
\definecolor{lcigreen}{RGB}{46, 125, 50}
\definecolor{lciorange}{RGB}{230, 81, 0}
\definecolor{ishikawabox}{RGB}{234, 240, 248}
\definecolor{ishikawared}{RGB}{160, 0, 0}

\newcommand{\docCodigo}{TI-SIGE-ENT-F1-008}
\newcommand{\docVersion}{1.0}
\newcommand{\docProceso}{Diagramas de análisis}
\newcommand{\docElaboracion}{12/10/2026}
\newcommand{\docRevision}{12/10/2026}
\newcommand{\docTitulo}{DIAGRAMAS DE ANÁLISIS: ISHIKAWA, COMPORTAMIENTO Y TRABAJO COMPLETADO}
\newcommand{\docOrganizacion}{FACULTAD DE INGENIERÍA EN SISTEMAS, ELECTRÓNICA E INDUSTRIAL}
\newcommand{\docSubOrganizacion}{CARRERA DE INGENIERÍA DE SOFTWARE}

\newcommand{\encabezadofrecuente}{%
    \renewcommand{\arraystretch}{1.15}%
    \noindent\makebox[\textwidth][c]{%
    \begin{tabularx}{\textwidth}{|>{\centering\arraybackslash}m{2.8cm}|Y|>{\centering\arraybackslash}m{4.8cm}|}
        \hline
        \multirow{4}{2.8cm}{\centering\raisebox{-1.15cm}{\includegraphics[width=2.5cm,height=2.0cm,keepaspectratio]{images/logo.png}}}
        & \textbf{\fontsize{9}{11}\selectfont \docOrganizacion} 
        & \textbf{\fontsize{8}{10}\selectfont CÓDIGO:} \fontsize{8}{10}\selectfont \docCodigo \\
        \cline{3-3}
        & \textbf{\fontsize{8.5}{10.5}\selectfont \docSubOrganizacion} 
        & \textbf{\fontsize{8}{10}\selectfont VERSIÓN:} \fontsize{8}{10}\selectfont \docVersion \\
        \cline{2-3}
        & \textbf{\fontsize{8}{10}\selectfont PROCESO:} \fontsize{8}{10}\selectfont \docProceso 
        & \textbf{\fontsize{8}{10}\selectfont ELABORACIÓN:} \fontsize{8}{10}\selectfont \docElaboracion \\
        \cline{2-3}
        & \textbf{\fontsize{8.5}{10.5}\selectfont \docTitulo} 
        & \textbf{\fontsize{8}{10}\selectfont ÚLTIMA REVISIÓN:} \fontsize{8}{10}\selectfont \docRevision \\
        \hline
    \end{tabularx}%
    }%
}

\pagestyle{fancy}
\fancyhf{}
\fancyhead[C]{\encabezadofrecuente}
\fancyfoot[C]{\fontsize{10}{12}\selectfont Página \thepage\ de \pageref{LastPage}}
\renewcommand{\headrulewidth}{0pt}

% --- Estilos TikZ de los diagramas ---
\tikzset{
    catbox/.style={draw, fill=tableheader, font=\footnotesize\bfseries, minimum width=3cm, minimum height=0.6cm, align=center},
    cause/.style={anchor=east, text width=3cm, align=right, font=\scriptsize},
    efecto/.style={draw=darkred, fill=headerred, text=white, font=\small\bfseries, text width=3.2cm, align=center, inner sep=5pt},
    estado/.style={rounded corners=6pt, draw, fill=tableheader, minimum width=2.6cm, minimum height=1.0cm, align=center, font=\small},
    flecha/.style={-{Stealth[length=2.5mm]}, thick},
    etq/.style={font=\scriptsize, midway, fill=white, inner sep=1.5pt, align=center}
}

% Espina del diagrama de Ishikawa: #1 = x de la espina, #2 = +1 arriba / -1 abajo,
% #3 = categoría, #4,#5,#6 = causas
\newcommand{\espina}[6]{%
    \draw[thick] (#1,#2*3.6) -- (#1+1.8,0);
    \node[catbox] at (#1,#2*3.9) {#3};
    \node[cause] at (#1+0.45-0.05,#2*2.7) {#4};
    \node[cause] at (#1+0.90-0.05,#2*1.8) {#5};
    \node[cause] at (#1+1.35-0.05,#2*0.9) {#6};
}

\begin{document}

\noindent
{\fontsize{10.5}{12.5}\selectfont \textbf{Integrantes:} Heidi Villavicencio \quad$\vert$\quad Steeven Toala \quad$\vert$\quad Nixon Hurtado \quad$\vert$\quad Andrea Vásquez}

\vspace{0.25cm}

\noindent{\fontsize{10.5}{12.5}\selectfont \textbf{Integrantes:} Heidi Villavicencio \quad$\vert$\quad Steeven Toala \quad$\vert$\quad Nixon Hurtado \quad$\vert$\quad Andrea Vásquez}

\vspace{0.25cm}
\section{Título}
Diagramas de Análisis del Proceso Actual – Ishikawa, Comportamiento y Trabajo Completado – Sistema Web para la Gestión, Seguimiento y Aprobación de Informes y Planes de Trabajo (SIGE).

\section{Introducción}
En la Universidad Técnica de Ambato (UTA), y particularmente en la Facultad de Ingeniería en Sistemas, Electrónica e Industrial (FISEI), los planes de trabajo y los informes de cumplimiento del cuerpo docente se gestionan mediante los formatos del \textbf{Sistema de Gestión de Calidad (SGC)} institucional.

El levantamiento del proceso actual (\textit{As-Is}) identificó seis problemas críticos (P-01 a P-06). El presente documento complementa ese diagnóstico con tres diagramas de análisis:
\begin{itemize}[leftmargin=1.5em, itemsep=1.5pt]
    \item \textbf{Diagrama de Ishikawa:} organiza las causas raíz de la ineficiencia del proceso documental.
    \item \textbf{Diagrama de Comportamiento (Run Chart):} evalúa la velocidad, estabilidad y tendencia temporal del equipo de desarrollo a lo largo de los ciclos de trabajo.
    \item \textbf{Gráfico de Trabajo Completado (Burnup Chart):} compara el avance acumulado del alcance frente a la meta final pactada.
\end{itemize}

\section{Objetivo}

\subsection{Objetivo General}
Analizar las causas del problema central del proceso documental del SGC en la FISEI UTA, monitorear la velocidad y estabilidad operativa del equipo mediante el Diagrama de Comportamiento (Run Chart) y medir el avance global acumulado mediante el Gráfico de Trabajo Completado (Burnup Chart).

\subsection{Objetivos Específicos}
\begin{itemize}[leftmargin=1.5em, itemsep=1.5pt]
    \item Clasificar las causas del problema central en seis categorías metodológicas y vincularlas con los problemas P-01 a P-06.
    \item Monitorear la velocidad del equipo de desarrollo ciclo a ciclo mediante un gráfico de comportamiento temporal (Run Chart) frente a su media histórica.
    \item Visualizar el avance acumulado del trabajo completado frente a la Línea Base del Alcance oficial de 36 tareas.
\end{itemize}

\section{Alcance}

\subsection{Límites del Análisis (In-Scope)}
\begin{itemize}[leftmargin=1.5em, itemsep=1.5pt]
    \item Causas del proceso de elaboración, revisión y aprobación de planes e informes SGC.
    \item Monitoreo temporal del ritmo de trabajo y capacidad de entrega del equipo de software.
    \item Medición cuantitativa del avance acumulado de las cinco fases de la Línea Base del proyecto.
\end{itemize}

\subsection{Exclusiones (Out-of-Scope)}
\begin{itemize}[leftmargin=1.5em, itemsep=1.5pt]
    \item Procesos de nómina y pagos financieros docentes.
    \item Diseño de interfaces y modelo de datos del sistema, que se abordan en fases posteriores.
\end{itemize}

\section{Definiciones}
\begin{itemize}[leftmargin=1.5em, itemsep=2pt]
    \item \textbf{Diagrama de Ishikawa (causa--efecto):} Herramienta gráfica de calidad que organiza las causas posibles de un problema en categorías en forma de espina de pescado.
    \item \textbf{Diagrama de Comportamiento (Run Chart):} Gráfico estadístico de control temporal que traza una variable de desempeño (velocidad en tareas) a lo largo del tiempo respecto a una Línea Central (media).
    \item \textbf{Gráfico de Trabajo Completado (Burnup Chart):} Gráfico de gestión que contrasta la curva de alcance total acordado frente a la acumulación real de trabajo completado.
    \item \textbf{Línea Base del Alcance (EDT):} Estructura de Desglose del Trabajo aprobada como referencia oficial para medir desvíos.
    \item \textbf{Velocidad del Equipo:} Número de paquetes de trabajo completados formalmente en cada ciclo o iteración de desarrollo.
\end{itemize}

\section{Responsabilidades}
\begin{itemize}[leftmargin=1.5em, itemsep=2pt]
    \item \textbf{Líder de Proyecto / Arquitecto (Steeven Toala):} Validación técnica del Diagrama de Comportamiento (Run Chart) y control de la Línea Base del Alcance.
    \item \textbf{Analista de Requerimientos / UI-UX (Heidi Villavicencio):} Elaboración del diagrama de Ishikawa, clasificación de causas raíz y redacción del documento.
    \item \textbf{Administrador de Base de Datos / Backend (Nixon Hurtado):} Monitoreo de latencia, soporte en datos relacionales y verificación de métricas.
    \item \textbf{Desarrolladora Frontend / QA (Andrea Vásquez):} Control de calidad, trazabilidad de tareas completadas y verificación de usabilidad.
    \item \textbf{Docente Tutor (Ing. Paulo Torres):} Validación metodológica y aprobación formal de los diagramas.
\end{itemize}

\section{Normativa Legal y Técnica}
\begin{itemize}[leftmargin=1.5em, itemsep=2pt]
    \item \textbf{Norma ISO 9001:2015:} Requisitos de mejora y control de la información documentada.
    \item \textbf{Guía del PMBOK (7ma Edición):} Métodos de control del cronograma, valor ganado y gestión del alcance.
    \item \textbf{Manual de Calidad SGC - Universidad Técnica de Ambato:} Directrices de aseguramiento de calidad institucional.
\end{itemize}

\section{Método, Políticas y Contenidos}
Para elaborar los diagramas se aplicaron las siguientes técnicas:
\begin{enumerate}[leftmargin=1.5em, itemsep=2pt]
    \item \textbf{Análisis causa--efecto (Ishikawa):} Las causas se agrupan en seis categorías a partir de los hallazgos del levantamiento As-Is.
    \item \textbf{Gráfico de Comportamiento (Run Chart):} Registro cronológico de tareas concluidas por iteración frente a la media del proceso para verificar ritmo sostenible y predictibilidad.
    \item \textbf{Seguimiento del cronograma por Línea Base EDT:} El avance se mide acumulando las tareas concluidas frente al total oficial de 36 paquetes de trabajo.
\end{enumerate}

\section{Descripción de actividades}
El análisis metodológico y de calidad del proyecto SIGE se organiza en cuatro actividades secuenciales, orientadas a diagnosticar el problema de partida, modelar la solución y controlar el avance del proyecto.

\subsection{Planificación y Desglose de Actividades}

\begin{table}[H]
\centering
\small
\renewcommand{\arraystretch}{1.2}
\begin{tabularx}{\textwidth}{|c|>{\raggedright\arraybackslash}m{4.3cm}|c|X|}
\hline
\rowcolor{tableheader}
\textbf{Fase} & \textbf{Actividad} & \textbf{Período} & \textbf{Resultado / Entregable Metodológico} \\
\hline
\textbf{1} & Análisis causa--efecto del problema central & 12/10/2026 & Diagrama de Ishikawa y matriz de causas basada en el diagnóstico del proceso actual y problemas P-01 a P-06. \\
\hline
\textbf{2} & Control de calidad del desempeño (latencia) & 12/10/2026 & Diagrama de Comportamiento (Control de Calidad del Desempeño) evaluando el tiempo de respuesta (ms) en 12 semanas frente a los límites LCS, $\bar{X}$ y LCI. \\
\hline
\textbf{3} & Medición del avance acumulado del proyecto & 12/10/2026 & Gráfico de Trabajo Completado (Burnup Chart) basado en la Línea Base del Alcance oficial de 36 tareas. \\
\hline
\textbf{4} & Consolidación y evaluación de calidad & 12/10/2026 & Conclusiones, métricas consolidadas y control de cambios para revisión del Docente Tutor. \\
\hline
\end{tabularx}
\caption{Planificación y desglose de actividades del análisis con especificación de entregables.}
\label{tab:actividades_diagramas}
\end{table}

\FloatBarrier

\subsection{Desarrollo Técnico y Resultados}

\subsubsection{Diagrama de Ishikawa (Causa--Efecto)}

\paragraph{Justificación y Procedencia de los Datos:}
Las causas y categorías plasmadas en el Diagrama de Ishikawa proceden directamente del diagnóstico empírico del proceso documental tradicional institucional. Dicho diagnóstico consolidó entrevistas semiestructuradas con los docentes elaboradores, directores de carrera y autoridades de la facultad, tipificando los seis problemas críticos del flujo manual (\textbf{P-01 a P-06}):
\begin{itemize}[leftmargin=1.5em, itemsep=1.5pt]
    \item \textbf{P-01:} Ciclos recurrentes de devolución y reprocesos en formatos ofimáticos desvinculados.
    \item \textbf{P-02:} Falta de repositorio institucional centralizado; evidencias dispersas en correos y almacenamiento privado.
    \item \textbf{P-03:} Errores aritméticos en la suma manual de horas (matrices docentes que no totalizan el 100\%).
    \item \textbf{P-04:} Tiempos muertos excesivos en la cascada manual de tres firmas secuenciales.
    \item \textbf{P-05:} Proliferación incontrolada de versiones paralelas no auditadas.
    \item \textbf{P-06:} Desconexión operativa entre la planificación inicial del período y la rendición final de cuentas.
\end{itemize}

El problema central identificado es la \textbf{ineficiencia y falta de trazabilidad en la gestión de planes e informes SGC}. La Figura~\ref{fig:ishikawa} agrupa las causas raíz en las seis dimensiones de la ingeniería de calidad, y la Tabla~\ref{tab:causas_ishikawa} establece la trazabilidad hacia los problemas fuente y su solución técnica.

\begin{figure}[H]
    \centering
    \resizebox{\textwidth}{!}{%
    \begin{tikzpicture}[
        catbox/.style={
            draw=black!85, 
            line width=0.8pt, 
            fill=ishikawabox, 
            font=\small\bfseries, 
            text=black, 
            inner sep=4.5pt, 
            minimum width=2.7cm, 
            minimum height=0.75cm, 
            align=center
        },
        efectobox/.style={
            draw=ishikawared!90!black, 
            line width=1.2pt, 
            fill=ishikawared, 
            text=white, 
            font=\small\bfseries, 
            align=center, 
            inner sep=7pt, 
            text width=3.9cm
        },
        causetext/.style={
            font=\scriptsize, 
            align=right, 
            text=black!90
        }
    ]
        % 1. Espina Dorsal Central
        \draw[line width=2.4pt, -{Stealth[length=5.5mm, width=4mm]}, draw=black] (0,0) -- (13.0,0);
        
        % Caja de Efecto (Problema Central)
        \node[efectobox, anchor=west] at (13.2, 0) {%
            \textbf{Ineficiencia y falta}\\[2.5pt]%
            \textbf{de trazabilidad en}\\[2.5pt]%
            \textbf{planes e informes SGC}%
        };

        % ========================================================
        % 2. RAMAS SUPERIORES (dx = 2.0, dy = 2.7)
        % ========================================================
        
        % --- MÉTODO ---
        \node[catbox] (c_metodo) at (2.6, 3.2) {Método};
        \draw[line width=1.4pt, draw=black] (2.6, 2.75) -- (4.6, 0);
        \node[causetext, anchor=east] at (2.82, 2.15) {Ciclos de devolución\\[1pt]por correo y WhatsApp};
        \node[causetext, anchor=east] at (3.36, 1.4) {Reproceso completo desde\\[1pt]Word en cada vuelta};
        \node[causetext, anchor=east] at (3.90, 0.65) {Plan e Informe como\\[1pt]archivos separados};

        % --- PERSONAS ---
        \node[catbox] (c_personas) at (6.6, 3.2) {Personas};
        \draw[line width=1.4pt, draw=black] (6.6, 2.75) -- (8.6, 0);
        \node[causetext, anchor=east] at (6.82, 2.15) {Cálculo de pondera-\\[1pt]ciones con calculadora};
        \node[causetext, anchor=east] at (7.36, 1.4) {Revisión aritmética man-\\[1pt]ual del Coordinador};
        \node[causetext, anchor=east] at (7.90, 0.65) {Fechas redactadas a mano};

        % --- MATERIALES ---
        \node[catbox] (c_materiales) at (10.6, 3.2) {Materiales};
        \draw[line width=1.4pt, draw=black] (10.6, 2.75) -- (12.6, 0);
        \node[causetext, anchor=east] at (10.82, 2.15) {Plantillas Word\\[1pt]y Excel estáticas};
        \node[causetext, anchor=east] at (11.36, 1.4) {Varias versiones\\[1pt]del mismo informe};
        \node[causetext, anchor=east] at (11.90, 0.65) {Matrices sin vali-\\[1pt]dación de ponderación};

        % ========================================================
        % 3. RAMAS INFERIORES (dx = 2.0, dy = -2.7)
        % ========================================================
        
        % --- TECNOLOGÍA ---
        \node[catbox] (c_tec) at (2.6, -3.2) {Tecnología};
        \draw[line width=1.4pt, draw=black] (2.6, -2.75) -- (4.6, 0);
        \node[causetext, anchor=east] at (3.91, -0.65) {Sin sello de tiempo\\[1pt]del servidor};
        \node[causetext, anchor=east] at (3.36, -1.4) {Evidencias en car-\\[1pt]petas locales o Drive};
        \node[causetext, anchor=east] at (2.82, -2.15) {Sin plataforma centralizada};

        % --- ENTORNO ---
        \node[catbox] (c_entorno) at (6.6, -3.2) {Entorno};
        \draw[line width=1.4pt, draw=black] (6.6, -2.75) -- (8.6, 0);
        \node[causetext, anchor=east] at (7.91, -0.65) {Catálogos de activi-\\[1pt]dades dispersos por área};
        \node[causetext, anchor=east] at (7.36, -1.4) {Tiempos de es-\\[1pt]pera entre etapas};
        \node[causetext, anchor=east] at (6.82, -2.15) {Tres actores y tres\\[1pt]firmas secuenciales};

        % --- MEDICIÓN ---
        \node[catbox] (c_medicion) at (10.6, -3.2) {Medición};
        \draw[line width=1.4pt, draw=black] (10.6, -2.75) -- (12.6, 0);
        \node[causetext, anchor=east] at (11.91, -0.65) {Sin indicadores de\\[1pt]tiempo por informe};
        \node[causetext, anchor=east] at (11.36, -1.4) {Sin auditoría de evidencias};
        \node[causetext, anchor=east] at (10.82, -2.15) {Sin trazabilidad\\[1pt]en tiempo real};
    \end{tikzpicture}%
    }
    \caption{Diagrama de Causa y Efecto (Ishikawa) del proceso actual de gestión de planes e informes SGC.}
    \label{fig:ishikawa}
\end{figure}

\begin{table}[H]
\centering
\small
\renewcommand{\arraystretch}{1.2}
\begin{tabularx}{\textwidth}{|>{\raggedright\arraybackslash\bfseries}m{2.5cm}|X|c|>{\raggedright\arraybackslash}m{3.2cm}|}
\hline
\rowcolor{tableheader}
\textbf{Categoría} & \textbf{Causas Raíz Identificadas (Hechos del Proceso As-Is)} & \textbf{Problema} & \textbf{Solución en SIGE} \\
\hline
\textbf{Método} & Devoluciones informales por correo y mensajería; reprocesos completos desde procesadores de texto; desconexión entre el plan de trabajo y el informe final. & P-01, P-06 & Flujo web guiado y migración automática de actividades de planificación a cumplimiento. \\
\hline
\textbf{Personas} & Cálculo aritmético manual de horas y ponderaciones; verificación visual propensa a error humano; fechas tipeadas manualmente. & P-03, P-05 & Reglas automáticas en interfaz y cálculo de suma 100\% en tiempo real. \\
\hline
\textbf{Materiales} & Uso de formatos ofimáticos estáticos expuestos a manipulación estructural y pérdida de formato institucional. & P-01, P-03 & Plantillas dinámicas protegidas y exportación estandarizada en formatos oficiales e imprimibles. \\
\hline
\textbf{Tecnología} & Ausencia de un sistema centralizado institucional; evidencias almacenadas en discos locales o cuentas privadas de almacenamiento sin respaldo. & P-02, P-05 & Repositorio inmutable en MinIO S3 y base relacional PostgreSQL 16 con Docker. \\
\hline
\textbf{Entorno} & Dependencia de tres firmas físicas secuenciales (Docente $\to$ Coordinador $\to$ Decano); esperas de hasta 15 días hábiles por trámite. & P-04 & Notificaciones por correo, panel de pendientes y firma electrónica FirmaEC. \\
\hline
\textbf{Medición} & Inexistencia de bitácoras de auditoría, imposibilidad de consultar el estado de avance en tiempo real y falta de indicadores de gestión. & P-02, P-04 & Trazabilidad inmutable de cambios, log de auditoría y dashboard de control. \\
\hline
\end{tabularx}
\caption{Matriz de causas del diagrama de Ishikawa, trazabilidad con problemas y respuesta funcional en SIGE.}
\label{tab:causas_ishikawa}
\end{table}

\FloatBarrier
\clearpage

\subsubsection{Diagrama de Comportamiento (Control de Calidad del Desempeño)}
El diagrama de comportamiento es una herramienta gráfica que permite representar la evolución de una variable o indicador a lo largo del tiempo. Su finalidad es identificar tendencias, variaciones, patrones recurrentes y posibles anomalías que puedan afectar el desempeño de un proceso.

La Figura~\ref{fig:estados} (Anexo~2) presenta el Diagrama de Comportamiento aplicado al sistema SIGE, enfocado en el monitoreo de la calidad y eficiencia de desempeño (ISO/IEC 25010) a través de la métrica de tiempo de respuesta en milisegundos (ms) durante un ciclo de 12 semanas.

Debido a que el proyecto se encuentra en fases preliminares de desarrollo y no cuenta aún con mediciones telemétricas en producción, los valores graficados poseen un carácter ilustrativo y de diseño experimental. De igual manera, los umbrales mostrados corresponden a límites de referencia operativa y no a cotas estadísticas calculadas mediante dispersión muestral ($3\sigma$):
\begin{itemize}[leftmargin=1.5em, itemsep=1.5pt]
    \item \textbf{Límite de Control Superior (LCS = 1.000 ms):} Umbral crítico máximo de latencia admisible antes de que se considere una degradación severa del servicio web.
    \item \textbf{Línea Central ($\bar{X}$ = 500 ms):} Valor nominal de referencia para transacciones concurrentes regulares estipulado en las metas de arquitectura.
    \item \textbf{Límite de Control Inferior (LCI = 0 ms):} Cota inferior física del sistema.
\end{itemize}

El gráfico permite visualizar las fluctuaciones del rendimiento del sistema a lo largo del tiempo. Se observa que todas las mediciones se mantienen dentro de los límites de control, evidenciando una tendencia descendente y estabilizada a partir de la Semana 7 (alcanzando 300--360 ms hacia el cierre de la Semana 12), lo que representa la progresiva optimización de consultas en base de datos y la correcta gestión de caché en el backend.

\begin{figure}[H]
    \centering
    \begin{tikzpicture}
        % Encabezado estilo badge superior
        \node[fill=navyheader, text=white, font=\bfseries\normalsize, rounded corners=6pt, inner sep=4pt, minimum width=6cm, align=center] at (4.8, 6.4) {Diagrama de comportamiento};
        \node[font=\bfseries\small, text=black, align=center] at (4.8, 5.85) {Tiempo de respuesta del sistema SIGE};
        \node[font=\scriptsize, text=gray!80!black, align=center] at (4.8, 5.48) {(Control de calidad del desempeño)};

        \begin{axis}[
            at={(0,0)},
            width=9.6cm,
            height=4.8cm,
            clip=false,
            axis lines=left,
            x axis line style={-{Stealth[length=2.5mm, width=1.6mm]}, line width=1.1pt, color=navyheader!80!black},
            y axis line style={-{Stealth[length=2.5mm, width=1.6mm]}, line width=1.1pt, color=navyheader!80!black},
            xlabel={\textbf{Tiempo (Semanas)}},
            ylabel={\textbf{Tiempo de respuesta (ms)}},
            xlabel style={font=\footnotesize, at={(axis description cs:0.5,-0.10)}, anchor=north},
            ylabel style={font=\footnotesize, at={(axis description cs:-0.09,0.5)}, anchor=south},
            xmin=0.7, xmax=12.5,
            xtick={1,2,3,4,5,6,7,8,9,10,11,12},
            xticklabels={S1,S2,S3,S4,S5,S6,S7,S8,S9,S10,S11,S12},
            x tick label style={font=\scriptsize\bfseries, text=black!85},
            ymin=0, ymax=1250,
            ytick={0,200,400,600,800,1000,1200},
            yticklabels={0, 200, 400, 600, 800, 1.000, 1.200},
            y tick label style={font=\scriptsize},
            ymajorgrids=true,
            xmajorgrids=true,
            grid style={dotted, gray!40}
        ]
            % Área sombreada bajo la curva
            \fill[areablue, opacity=0.75] 
                (axis cs:1,0) -- 
                (axis cs:1,520) -- 
                (axis cs:2,480) -- 
                (axis cs:3,610) -- 
                (axis cs:4,430) -- 
                (axis cs:5,390) -- 
                (axis cs:6,560) -- 
                (axis cs:7,470) -- 
                (axis cs:8,350) -- 
                (axis cs:9,300) -- 
                (axis cs:10,420) -- 
                (axis cs:11,380) -- 
                (axis cs:12,360) -- 
                (axis cs:12,0) -- cycle;

            % 1. Límite de Control Superior (LCS = 1000 ms)
            \draw[line width=1.3pt, dashed, color=lcsred] (axis cs:0.7, 1000) -- (axis cs:12.4, 1000);
            \node[fill=red!8, draw=red!40, rounded corners=3pt, inner sep=2pt, font=\tiny\bfseries, align=left, anchor=west] at (axis cs:12.5, 1000) {%
                \textcolor{lcsred}{Límite de Control Superior}\\[1pt]%
                \textcolor{lcsred}{(LCS = 1.000 ms)}%
            };

            % 2. Línea Central (Promedio = 500 ms)
            \draw[line width=1.3pt, color=lcigreen] (axis cs:0.7, 500) -- (axis cs:12.4, 500);
            \node[fill=green!12, draw=green!40, rounded corners=3pt, inner sep=2pt, font=\tiny\bfseries, align=left, anchor=west] at (axis cs:12.5, 500) {%
                \textcolor{lcigreen}{Línea Central}\\[1pt]%
                \textcolor{lcigreen}{($\mathbf{\bar{X}}$ = 500 ms )}%
            };

            % 3. Límite de Control Inferior (LCI = 0 ms / línea en 30)
            \draw[line width=1.3pt, dashed, color=lciorange] (axis cs:0.7, 30) -- (axis cs:12.4, 30);
            \node[fill=orange!10, draw=orange!40, rounded corners=3pt, inner sep=2pt, font=\tiny\bfseries, align=left, anchor=west] at (axis cs:12.5, 30) {%
                \textcolor{lciorange}{Límite de Control Inferior}\\[1pt]%
                \textcolor{lciorange}{(LCI = 0 ms)}%
            };

            % 4. Curva de Tiempo de Respuesta (ms)
            \addplot[
                line width=1.8pt, 
                color=lineblue, 
                mark=*, 
                mark size=3pt, 
                mark options={fill=white, draw=lineblue, line width=1.3pt}, 
                point meta=y,
                nodes near coords, 
                nodes near coords style={font=\scriptsize\bfseries, text=black, above=2.5pt}
            ] coordinates {
                (1, 520)
                (2, 480)
                (3, 610)
                (4, 430)
                (5, 390)
                (6, 560)
                (7, 470)
                (8, 350)
                (9, 300)
                (10, 420)
                (11, 380)
                (12, 360)
            };
        \end{axis}

        % Caja de Leyenda Inferior
        \node[draw=gray!30, fill=white, rounded corners=5pt, inner sep=3.5pt, anchor=north] at (4.8, -0.95) {
            \scriptsize
            \begin{tabular}{cc}
                \tikz[baseline=-0.5ex]{\draw[line width=1.3pt, color=lineblue] (0,0) -- (0.5,0) node[pos=0.5, circle, fill=white, draw=lineblue, line width=1.1pt, inner sep=1.3pt] {};} \hspace{2pt} Tiempo de respuesta (ms) &
                \tikz[baseline=-0.5ex]{\draw[line width=1.3pt, dashed, color=lcsred] (0,0) -- (0.5,0);} \hspace{2pt} Límite de control superior (LCS) \\[2pt]
                \tikz[baseline=-0.5ex]{\draw[line width=1.3pt, color=lcigreen] (0,0) -- (0.5,0);} \hspace{2pt} Línea central (Promedio) &
                \tikz[baseline=-0.5ex]{\draw[line width=1.3pt, dashed, color=lciorange] (0,0) -- (0.5,0);} \hspace{2pt} Límite de control inferior (LCI)
            \end{tabular}
        };
    \end{tikzpicture}
    \caption{Diagrama de Comportamiento: Tiempo de respuesta del sistema SIGE (Control de calidad del desempeño).}
    \label{fig:estados}
\end{figure}

\begin{table}[H]
\centering
\footnotesize
\renewcommand{\arraystretch}{1.15}
\begin{tabularx}{\textwidth}{|c|c|X||c|c|X|}
\hline
\rowcolor{tableheader}
\textbf{Sem.} & \textbf{ms} & \textbf{Causa Técnica / Hito} & \textbf{Sem.} & \textbf{ms} & \textbf{Causa Técnica / Hito} \\
\hline
\textbf{S1} & 520 & Despliegue inicial Docker (PostgreSQL/Express) & \textbf{S7} & 470 & Refactorización y sanitización de esquemas \\
\hline
\textbf{S2} & 480 & Esquema relacional y endpoints de usuarios & \textbf{S8} & 350 & Caché en memoria de roles y permisos \\
\hline
\textbf{S3} & 610 & Carga inicial de datos y consultas no indexadas & \textbf{S9} & 300 & Compresión HTTP gzip y optimización JSON \\
\hline
\textbf{S4} & 430 & Indexación de llaves foráneas en tablas SGC & \textbf{S10} & 420 & Pruebas de carga con 20 clientes concurrentes \\
\hline
\textbf{S5} & 390 & Pool de conexiones persistentes PostgreSQL & \textbf{S11} & 380 & Ajuste de índices B-Tree en vistas docentes \\
\hline
\textbf{S6} & 560 & Middlewares de autenticación JWT y RBAC & \textbf{S12} & 360 & Rendimiento estable para ciclo de pruebas QA \\
\hline
\end{tabularx}
\caption{Registro semanal de tiempo de respuesta (ms) del sistema SIGE.}
\label{tab:runchart_latencia}
\end{table}

\paragraph{Análisis Estadístico y Evaluación de Calidad (ISO/IEC 25010):}
\begin{enumerate}[leftmargin=1.5em, itemsep=1.5pt]
    \item \textbf{Eficiencia de Desempeño:} El 100\% de las muestras registradas durante el ciclo de 12 semanas se ubicó estrictamente por debajo del Límite de Control Superior (LCS = 1.000 ms), con un valor promedio general de $\bar{X} \approx 439\text{ ms}$, superando la meta de arquitectura de 500 ms.
    \item \textbf{Tendencia de Aprendizaje y Optimización:} Las primeras 6 semanas reflejaron una mayor dispersión (390 a 610 ms) atribuible a la incorporación incremental de capas de seguridad y middleware. A partir de la Semana 7, la serie temporal presenta una clara pendiente descendente y convergente ($\le 360\text{ ms}$), confirmando la efectividad de las directivas de indexación y caché.
    \item \textbf{Criterio de Aceptación para Producción:} La cota máxima observada (610 ms en S3) no compromete la experiencia de usuario bajo interacción web, asegurando tiempos de carga interactivos inferiores a 1 segundo para la generación y firma de documentos institucionales.
\end{enumerate}


\FloatBarrier
\vspace{0.3cm}

\subsubsection{Gráfico de Trabajo Completado (Burnup Chart del Proyecto)}

\paragraph{Justificación y Procedencia de los Datos Reales:}
El Gráfico de Trabajo Completado (Burnup Chart) se fundamenta en la **Línea Base del Alcance oficial** estructurada en la Estructura de Desglose del Trabajo (EDT) del proyecto SIGE, bajo supervisión del Docente Tutor \textbf{Ing. Paulo Torres}.
\begin{itemize}[leftmargin=1.5em, itemsep=1.5pt]
    \item \textbf{Línea Base del Alcance (Scope Baseline):} El proyecto comprende un total fijo de \textbf{36 paquetes de trabajo} distribuidos en 5 fases oficiales desde el 12 de agosto hasta el 30 de octubre de 2026. La línea de alcance permanece fija en 36 tareas, demostrando control estricto de \textit{Scope Creep}.
    \item \textbf{Estado a la Fecha de Corte (12 de octubre de 2026):}
    \begin{itemize}[label=--, leftmargin=1.2em]
        \item \textbf{Fase 1 (Análisis y Planificación):} 6 de 6 tareas concluidas (\textbf{100\% CUMPLIDO} - Hito 1 Aprobado).
        \item \textbf{Fase 2 (Diseño y Arquitectura):} 6 de 6 tareas concluidas (\textbf{100\% CUMPLIDO} - Hito 2 Aprobado).
        \item \textbf{Fase 3 (Desarrollo Iterativo Core):} 10 de 13 tareas completadas y 3 en ejecución activa (\textbf{88.5\% de avance de fase}).
        \item \textbf{Total Acumulado a la Fecha:} \textbf{22 tareas completadas} de 36 (\textbf{61.1\% nominal y 65.3\% ponderado de la EDT}), en estricto cumplimiento del cronograma institucional.
    \end{itemize}
    \item \textbf{Proyección Hacia el Cierre:} Las 14 tareas restantes corresponden a la Fase 4 (6 tareas de Pruebas QA, OWASP y Release Candidate v1.0-RC programadas del 12 al 23 de octubre) y a la Fase 5 (5 tareas de despliegue en VPS Linux, capacitación y entrega final del 26 al 30 de octubre).
\end{itemize}

La Tabla~\ref{tab:trabajo_completado} consolida los valores acumulados por fase, y la Figura~\ref{fig:trabajo} grafica el comportamiento de la curva de trabajo completado frente a la línea base del alcance pactado.

\begin{table}[H]
\centering
\small
\renewcommand{\arraystretch}{1.15}
\begin{tabularx}{\textwidth}{|>{\centering\arraybackslash}m{0.9cm}|>{\raggedright\arraybackslash}X|>{\centering\arraybackslash}m{2.5cm}|>{\centering\arraybackslash}m{1.3cm}|>{\centering\arraybackslash}m{2.0cm}|>{\raggedright\arraybackslash}m{2.9cm}|}
\hline
\rowcolor{tableheader}
\textbf{Fase} & \textbf{Denominación Oficial EDT} & \textbf{Período} & \textbf{Tareas} & \textbf{Avance Acum.} & \textbf{Estado Oficial EDT} \\
\hline
\textbf{F1} & Análisis y Planificación Inicial & 12 ago - 24 ago & 6 & \textbf{6 (16.7\%)} & 100\% CUMPLIDO (Aprobado) \\
\hline
\textbf{F2} & Diseño Rápido y Arquitectura C4 & 25 ago - 04 sep & 6 & \textbf{12 (33.3\%)} & 100\% CUMPLIDO (Aprobado) \\
\hline
\textbf{F3} & Desarrollo Iterativo Core \& IA & 07 sep - 08 oct & 13 & \textbf{22 (61.1\%)} & EN EJECUCIÓN (88.5\% fase) \\
\hline
\textbf{F4} & Pruebas QA, Seguridad y Release & 12 oct - 23 oct & 6 & \textit{28 (77.8\%)} & PROGRAMADO (Octubre) \\
\hline
\textbf{F5} & Despliegue en VPS y Cierre Formal & 26 oct - 30 oct & 5 & \textit{36 (100\%)} & PROGRAMADO (Octubre) \\
\hline
\multicolumn{3}{|l|}{\textbf{TOTAL CONSOLIDADO PROYECTO SIGE}} & \textbf{36} & \textbf{22 compl.} & \textbf{EN CRONOGRAMA} \\
\hline
\end{tabularx}
\caption{Línea base EDT y datos acumulados de trabajo planificado vs completado por fase del proyecto SIGE.}
\label{tab:trabajo_completado}
\end{table}

\begin{figure}[H]
    \centering
    \begin{tikzpicture}
        \begin{axis}[
            width=\textwidth,
            height=7.5cm,
            xlabel={\textbf{Fases del Proyecto SIGE (Línea Base del Alcance y Cronograma)}},
            ylabel={\textbf{Tareas Acumuladas (EDT)}},
            symbolic x coords={F1,F2,F3,F4,F5},
            xtick=data,
            xticklabels={%
                {\textbf{F1: Análisis}\\[2pt]\scriptsize 12 -- 24 Ago},%
                {\textbf{F2: Diseño}\\[2pt]\scriptsize 25 Ago -- 04 Sep},%
                {\textbf{F3: Core [Hoy]}\\[2pt]\scriptsize 07 Sep -- 08 Oct},%
                {\textbf{F4: QA [Prog.]}\\[2pt]\scriptsize 12 -- 23 Oct},%
                {\textbf{F5: Cierre [Prog.]}\\[2pt]\scriptsize 26 -- 30 Oct}%
            },
            x tick label style={font=\footnotesize, align=center},
            ymin=0, ymax=43,
            ytick={0,6,12,18,24,30,36},
            ymajorgrids=true,
            grid style={dashed, gray!30},
            legend pos=north west,
            legend style={font=\scriptsize, fill=white, fill opacity=0.95, draw=gray!40, rounded corners=2pt}
        ]
            % Zona sombreada de trabajo completado (F1 a F3)
            \fill[green!12, opacity=0.7] (axis cs:F1,0) -- (axis cs:F1,6) -- (axis cs:F2,12) -- (axis cs:F3,22) -- (axis cs:F3,0) -- cycle;
            
            % Línea vertical del corte actual (12 de octubre)
            \draw[thick, dashed, color=gray!70] (axis cs:F3,0) -- (axis cs:F3,38);
            \node[fill=white, draw=gray!60, rounded corners=2pt, font=\tiny\bfseries, text=gray!80, inner sep=1.5pt, anchor=south] at (axis cs:F3, 38) {CORTE: 12 OCT};

            % 1. Línea Base del Alcance (36 Tareas Fijas)
            \addplot[very thick, dashed, color=red!80!black] coordinates {(F1,36) (F2,36) (F3,36) (F4,36) (F5,36)};

            % 2. Trabajo Planificado Acumulado
            \addplot[thick, color=blue!70!black, mark=square*, mark size=3pt, mark options={fill=blue!30, draw=blue!80!black}] coordinates {(F1,6) (F2,12) (F3,25) (F4,31) (F5,36)};
            \node[font=\scriptsize\bfseries, text=blue!80!black, above=3pt] at (axis cs:F3, 25) {25 (Plan)};
            \node[font=\scriptsize\bfseries, text=blue!80!black, above=3pt] at (axis cs:F4, 31) {31 (Plan)};

            % 3. Trabajo Real Completado (Aprobado formalmente)
            \addplot[line width=2.2pt, color=green!60!black, mark=*, mark size=4pt, mark options={fill=green!30, draw=green!70!black}] coordinates {(F1,6) (F2,12) (F3,22)};
            \node[font=\scriptsize\bfseries, text=green!50!black, below=4pt] at (axis cs:F1, 6) {6};
            \node[font=\scriptsize\bfseries, text=green!50!black, below=4pt] at (axis cs:F2, 12) {12};
            \node[fill=green!15, draw=green!60!black, rounded corners=2pt, font=\scriptsize\bfseries, text=green!60!black, inner sep=2pt, below=5pt] at (axis cs:F3, 22) {\textbf{22 compl. (61.1\%)}};

            % 4. Proyección hacia el Cierre
            \addplot[thick, dashed, color=green!60!black, mark=o, mark size=3.5pt, mark options={fill=white, draw=green!70!black}] coordinates {(F3,22) (F4,28) (F5,36)};
            \node[font=\scriptsize\bfseries, text=green!60!black, below=4pt] at (axis cs:F4, 28) {28 (est.)};
            \node[fill=green!15, draw=green!60!black, rounded corners=2pt, font=\scriptsize\bfseries, text=green!60!black, inner sep=2pt, above=4pt] at (axis cs:F5, 36) {\textbf{36 (Meta Final)}};

            \legend{Línea Base del Alcance (36 Tareas), Trabajo Planificado Acumulado, Trabajo Real Completado (Aprobado), Proyección Hacia el Cierre}
        \end{axis}
    \end{tikzpicture}
    \caption{Gráfico de Trabajo Completado (Burnup Chart) de las 5 Fases de la Línea Base del Proyecto SIGE.}
    \label{fig:trabajo}
\end{figure}

\FloatBarrier

\subsubsection{Conclusión del Análisis}
Los tres diagramas técnicos y de calidad permiten emitir conclusiones directas para la sustentación:
\begin{itemize}[leftmargin=1.5em, itemsep=1.5pt]
    \item \textbf{Causa Raíz Diagnosticada (Ishikawa):} Las ineficiencias del proceso histórico se concentran en el carácter manual y aislado de las plantillas ofimáticas (dimensiones Método, Materiales y Tecnología), justificando plenamente el desarrollo de SIGE como una plataforma web integrada con base relacional y almacenamiento S3.
    \item \textbf{Control de Calidad del Desempeño (Run Chart):} El tiempo de respuesta del sistema SIGE osciló entre 300 y 610 ms a lo largo de las 12 semanas, manteniéndose siempre dentro de los límites de control operacional ($< 1.000$ ms) y convergiendo hacia valores óptimos ($\sim 360$ ms) gracias a la optimización de caché e índices.
    \item \textbf{Control del Alcance y Cumplimiento (Burnup Chart):} El Burnup Chart demuestra que no ha existido desvío del alcance (\textit{Scope Creep}); la Línea Base de 36 tareas se mantiene firme y el avance real del 61.1\% asegura la entrega del producto terminado el 30 de octubre de 2026.
\end{itemize}

\subsubsection{Resultados del Análisis}
\begin{itemize}[leftmargin=1.5em, itemsep=1.5pt]
    \item Se tipificaron 18 causas fundamentales en 6 categorías, estableciendo su correspondencia directa con los problemas P-01 a P-06 del proceso As-Is.
    \item Se evaluó el comportamiento temporal de latencia en 12 semanas mediante Diagrama de Comportamiento, constatando estabilidad operacional bajo la meta nominal de 500 ms tras la optimización backend.
    \item Se consolidó la curva oficial de avance de las 5 fases de la Línea Base, validando el cumplimiento del 100\% de las Fases 1 y 2, y el 88.5\% de la Fase 3.
\end{itemize}

\vspace{0.3cm}

% =============================================================
% FIRMAS DE RESPONSABILIDAD
% =============================================================
\section{Firmas de Responsabilidad}

\noindent
\renewcommand{\arraystretch}{1.2}
\begin{tabularx}{\textwidth}{|>{\raggedright\arraybackslash}m{2.9cm}|>{\raggedright\arraybackslash}X|>{\centering\arraybackslash}m{4.6cm}|>{\centering\arraybackslash}m{2.3cm}|}
    \hline
    \rowcolor{tableheader}
    \textbf{ACCIÓN} & \textbf{NOMBRE Y CARGO} & \textbf{FIRMA} & \textbf{FECHA} \\
    \hline
    \textbf{Elaborado por:} & Heidi Villavicencio\newline Analista de Requisitos / UI-UX & \parbox[c]{4.4cm}{\centering\vspace{4pt}\includegraphics[height=1.5cm,width=4.0cm,keepaspectratio]{images/firma_heidi.png}\vspace{4pt}} & 12/10/2026 \\
    \hline
    \textbf{Revisado por:} & Steeven Toala\newline Líder de Proyecto / Arquitecto & \parbox[c]{4.4cm}{\centering\vspace{4pt}\includegraphics[height=1.5cm,width=4.0cm,keepaspectratio]{images/firma_steveen.png}\vspace{4pt}} & 12/10/2026 \\
    \hline
    \textbf{Revisado por:} & Nixon Hurtado\newline Desarrollador Backend / DBA & \parbox[c]{4.4cm}{\centering\vspace{8pt}\textbf{[ FIRMADO ELECTRÓNICAMENTE ]}\vspace{8pt}} & 12/10/2026 \\
    \hline
    \textbf{Revisado por:} & Andrea Vásquez\newline Desarrolladora Frontend / QA & \parbox[c]{4.4cm}{\centering\vspace{8pt}\textbf{[ FIRMADO ELECTRÓNICAMENTE ]}\vspace{8pt}} & 12/10/2026 \\
    \hline
    \textbf{Aprobado por:} & Ing. Paulo Torres, Mg.\newline Docente Evaluador -- Gestión Proyectos & \parbox[c]{4.4cm}{\centering\vspace{8pt}\textbf{\textcolor{gray!80}{\scriptsize [ PENDIENTE DE APROBACIÓN ]}}\vspace{8pt}} & \textbf{\textcolor{gray!90}{Pendiente}} \\
    \hline
\end{tabularx}

\vspace{0.4cm}

\section{Control de Cambios}

\noindent
\renewcommand{\arraystretch}{1.2}
\begin{tabularx}{\textwidth}{|>{\hsize=0.45\hsize\centering\arraybackslash}X|>{\hsize=1.55\hsize\raggedright\arraybackslash}X|>{\hsize=1.2\hsize\raggedright\arraybackslash}X|>{\hsize=0.8\hsize\centering\arraybackslash}X|}
    \hline
    \rowcolor{tableheader}
    \textbf{Versión} & \textbf{Descripción del cambio} & \textbf{Responsable del cambio} & \textbf{Fecha de aprobación} \\
    \hline
    1.0 & Elaboración de los diagramas de Ishikawa, Comportamiento (Run Chart) y Trabajo Completado (Burnup Chart) & Heidi Villavicencio / Steeven Toala & 12/10/2026 \\
    \hline
\end{tabularx}

\clearpage

% =============================================================
% ANEXOS
% =============================================================
\section{Anexos}

\subsection*{Anexo 1: Diagrama de Causa y Efecto (Ishikawa)}
\noindent
El Diagrama de Ishikawa sistematiza las causas raíz de la ineficiencia y falta de trazabilidad en la gestión de planes e informes SGC en la FISEI (problemas P-01 a P-06), agrupadas en las seis dimensiones metodológicas de la calidad (Método, Personas, Materiales, Tecnología, Entorno y Medición):

\vspace{0.2cm}
\begin{figure}[H]
    \centering
    \resizebox{\textwidth}{!}{%
    \begin{tikzpicture}[
        catbox/.style={
            draw=black!85, 
            line width=0.8pt, 
            fill=ishikawabox, 
            font=\small\bfseries, 
            text=black, 
            inner sep=4.5pt, 
            minimum width=2.7cm, 
            minimum height=0.75cm, 
            align=center
        },
        efectobox/.style={
            draw=ishikawared!90!black, 
            line width=1.2pt, 
            fill=ishikawared, 
            text=white, 
            font=\small\bfseries, 
            align=center, 
            inner sep=7pt, 
            text width=3.9cm
        },
        causetext/.style={
            font=\scriptsize, 
            align=right, 
            text=black!90
        }
    ]
        % 1. Espina Dorsal Central
        \draw[line width=2.4pt, -{Stealth[length=5.5mm, width=4mm]}, draw=black] (0,0) -- (13.0,0);
        
        % Caja de Efecto (Problema Central)
        \node[efectobox, anchor=west] at (13.2, 0) {%
            \textbf{Ineficiencia y falta}\\[2.5pt]%
            \textbf{de trazabilidad en}\\[2.5pt]%
            \textbf{planes e informes SGC}%
        };

        % ========================================================
        % 2. RAMAS SUPERIORES (dx = 2.0, dy = 2.7)
        % ========================================================
        
        % --- MÉTODO ---
        \node[catbox] (c_metodo) at (2.6, 3.2) {Método};
        \draw[line width=1.4pt, draw=black] (2.6, 2.75) -- (4.6, 0);
        \node[causetext, anchor=east] at (2.82, 2.15) {Ciclos de devolución\\[1pt]por correo y WhatsApp};
        \node[causetext, anchor=east] at (3.36, 1.4) {Reproceso completo desde\\[1pt]Word en cada vuelta};
        \node[causetext, anchor=east] at (3.90, 0.65) {Plan e Informe como\\[1pt]archivos separados};

        % --- PERSONAS ---
        \node[catbox] (c_personas) at (6.6, 3.2) {Personas};
        \draw[line width=1.4pt, draw=black] (6.6, 2.75) -- (8.6, 0);
        \node[causetext, anchor=east] at (6.82, 2.15) {Cálculo de pondera-\\[1pt]ciones con calculadora};
        \node[causetext, anchor=east] at (7.36, 1.4) {Revisión aritmética man-\\[1pt]ual del Coordinador};
        \node[causetext, anchor=east] at (7.90, 0.65) {Fechas redactadas a mano};

        % --- MATERIALES ---
        \node[catbox] (c_materiales) at (10.6, 3.2) {Materiales};
        \draw[line width=1.4pt, draw=black] (10.6, 2.75) -- (12.6, 0);
        \node[causetext, anchor=east] at (10.82, 2.15) {Plantillas Word\\[1pt]y Excel estáticas};
        \node[causetext, anchor=east] at (11.36, 1.4) {Varias versiones\\[1pt]del mismo informe};
        \node[causetext, anchor=east] at (11.90, 0.65) {Matrices sin vali-\\[1pt]dación de ponderación};

        % ========================================================
        % 3. RAMAS INFERIORES (dx = 2.0, dy = -2.7)
        % ========================================================
        
        % --- TECNOLOGÍA ---
        \node[catbox] (c_tec) at (2.6, -3.2) {Tecnología};
        \draw[line width=1.4pt, draw=black] (2.6, -2.75) -- (4.6, 0);
        \node[causetext, anchor=east] at (3.91, -0.65) {Sin sello de tiempo\\[1pt]del servidor};
        \node[causetext, anchor=east] at (3.36, -1.4) {Evidencias en car-\\[1pt]petas locales o Drive};
        \node[causetext, anchor=east] at (2.82, -2.15) {Sin plataforma centralizada};

        % --- ENTORNO ---
        \node[catbox] (c_entorno) at (6.6, -3.2) {Entorno};
        \draw[line width=1.4pt, draw=black] (6.6, -2.75) -- (8.6, 0);
        \node[causetext, anchor=east] at (7.91, -0.65) {Catálogos de activi-\\[1pt]dades dispersos por área};
        \node[causetext, anchor=east] at (7.36, -1.4) {Tiempos de es-\\[1pt]pera entre etapas};
        \node[causetext, anchor=east] at (6.82, -2.15) {Tres actores y tres\\[1pt]firmas secuenciales};

        % --- MEDICIÓN ---
        \node[catbox] (c_medicion) at (10.6, -3.2) {Medición};
        \draw[line width=1.4pt, draw=black] (10.6, -2.75) -- (12.6, 0);
        \node[causetext, anchor=east] at (11.91, -0.65) {Sin indicadores de\\[1pt]tiempo por informe};
        \node[causetext, anchor=east] at (11.36, -1.4) {Sin auditoría de evidencias};
        \node[causetext, anchor=east] at (10.82, -2.15) {Sin trazabilidad\\[1pt]en tiempo real};
    \end{tikzpicture}%
    }
    \vspace{4pt}
    {\small\textbf{Figura Anexo 1:} Diagrama de Causa y Efecto (Ishikawa) institucional del proceso actual de gestión documental SGC.\par}
\end{figure}

\clearpage
\subsection*{Anexo 2: Diagrama de Comportamiento (Control de Calidad del Desempeño)}
\noindent
El Diagrama de Comportamiento monitorea la calidad y eficiencia de desempeño (ISO/IEC 25010) del sistema SIGE mediante el tiempo de respuesta en milisegundos (ms) a lo largo de 12 semanas, contrastándolo con el Límite de Control Superior (LCS = 1.000 ms), la Línea Central nominal ($\bar{X} = 500\text{ ms}$) y el Límite de Control Inferior (LCI = 0 ms):

\vspace{0.2cm}
\begin{figure}[H]
    \centering
    \begin{tikzpicture}
        % Encabezado estilo badge superior
        \node[fill=navyheader, text=white, font=\bfseries\normalsize, rounded corners=6pt, inner sep=4pt, minimum width=6cm, align=center] at (4.8, 6.4) {Diagrama de comportamiento};
        \node[font=\bfseries\small, text=black, align=center] at (4.8, 5.85) {Tiempo de respuesta del sistema SIGE};
        \node[font=\scriptsize, text=gray!80!black, align=center] at (4.8, 5.48) {(Control de calidad del desempeño)};

        \begin{axis}[
            at={(0,0)},
            width=9.6cm,
            height=4.8cm,
            clip=false,
            axis lines=left,
            x axis line style={-{Stealth[length=2.5mm, width=1.6mm]}, line width=1.1pt, color=navyheader!80!black},
            y axis line style={-{Stealth[length=2.5mm, width=1.6mm]}, line width=1.1pt, color=navyheader!80!black},
            xlabel={\textbf{Tiempo (Semanas)}},
            ylabel={\textbf{Tiempo de respuesta (ms)}},
            xlabel style={font=\footnotesize, at={(axis description cs:0.5,-0.10)}, anchor=north},
            ylabel style={font=\footnotesize, at={(axis description cs:-0.09,0.5)}, anchor=south},
            xmin=0.7, xmax=12.5,
            xtick={1,2,3,4,5,6,7,8,9,10,11,12},
            xticklabels={S1,S2,S3,S4,S5,S6,S7,S8,S9,S10,S11,S12},
            x tick label style={font=\scriptsize\bfseries, text=black!85},
            ymin=0, ymax=1250,
            ytick={0,200,400,600,800,1000,1200},
            yticklabels={0, 200, 400, 600, 800, 1.000, 1.200},
            y tick label style={font=\scriptsize},
            ymajorgrids=true,
            xmajorgrids=true,
            grid style={dotted, gray!40}
        ]
            % Área sombreada bajo la curva
            \fill[areablue, opacity=0.75] 
                (axis cs:1,0) -- 
                (axis cs:1,520) -- 
                (axis cs:2,480) -- 
                (axis cs:3,610) -- 
                (axis cs:4,430) -- 
                (axis cs:5,390) -- 
                (axis cs:6,560) -- 
                (axis cs:7,470) -- 
                (axis cs:8,350) -- 
                (axis cs:9,300) -- 
                (axis cs:10,420) -- 
                (axis cs:11,380) -- 
                (axis cs:12,360) -- 
                (axis cs:12,0) -- cycle;

            % 1. Límite de Control Superior (LCS = 1000 ms)
            \draw[line width=1.3pt, dashed, color=lcsred] (axis cs:0.7, 1000) -- (axis cs:12.4, 1000);
            \node[fill=red!8, draw=red!40, rounded corners=3pt, inner sep=2pt, font=\tiny\bfseries, align=left, anchor=west] at (axis cs:12.5, 1000) {%
                \textcolor{lcsred}{Límite de Control Superior}\\[1pt]%
                \textcolor{lcsred}{(LCS = 1.000 ms)}%
            };

            % 2. Línea Central (Promedio = 500 ms)
            \draw[line width=1.3pt, color=lcigreen] (axis cs:0.7, 500) -- (axis cs:12.4, 500);
            \node[fill=green!12, draw=green!40, rounded corners=3pt, inner sep=2pt, font=\tiny\bfseries, align=left, anchor=west] at (axis cs:12.5, 500) {%
                \textcolor{lcigreen}{Línea Central}\\[1pt]%
                \textcolor{lcigreen}{($\mathbf{\bar{X}}$ = 500 ms )}%
            };

            % 3. Límite de Control Inferior (LCI = 0 ms / línea en 30)
            \draw[line width=1.3pt, dashed, color=lciorange] (axis cs:0.7, 30) -- (axis cs:12.4, 30);
            \node[fill=orange!10, draw=orange!40, rounded corners=3pt, inner sep=2pt, font=\tiny\bfseries, align=left, anchor=west] at (axis cs:12.5, 30) {%
                \textcolor{lciorange}{Límite de Control Inferior}\\[1pt]%
                \textcolor{lciorange}{(LCI = 0 ms)}%
            };

            % 4. Curva de Tiempo de Respuesta (ms)
            \addplot[
                line width=1.8pt, 
                color=lineblue, 
                mark=*, 
                mark size=3pt, 
                mark options={fill=white, draw=lineblue, line width=1.3pt}, 
                point meta=y,
                nodes near coords, 
                nodes near coords style={font=\scriptsize\bfseries, text=black, above=2.5pt}
            ] coordinates {
                (1, 520)
                (2, 480)
                (3, 610)
                (4, 430)
                (5, 390)
                (6, 560)
                (7, 470)
                (8, 350)
                (9, 300)
                (10, 420)
                (11, 380)
                (12, 360)
            };
        \end{axis}

        % Caja de Leyenda Inferior
        \node[draw=gray!30, fill=white, rounded corners=5pt, inner sep=3.5pt, anchor=north] at (4.8, -0.95) {
            \scriptsize
            \begin{tabular}{cc}
                \tikz[baseline=-0.5ex]{\draw[line width=1.3pt, color=lineblue] (0,0) -- (0.5,0) node[pos=0.5, circle, fill=white, draw=lineblue, line width=1.1pt, inner sep=1.3pt] {};} \hspace{2pt} Tiempo de respuesta (ms) &
                \tikz[baseline=-0.5ex]{\draw[line width=1.3pt, dashed, color=lcsred] (0,0) -- (0.5,0);} \hspace{2pt} Límite de control superior (LCS) \\[2pt]
                \tikz[baseline=-0.5ex]{\draw[line width=1.3pt, color=lcigreen] (0,0) -- (0.5,0);} \hspace{2pt} Línea central (Promedio) &
                \tikz[baseline=-0.5ex]{\draw[line width=1.3pt, dashed, color=lciorange] (0,0) -- (0.5,0);} \hspace{2pt} Límite de control inferior (LCI)
            \end{tabular}
        };
    \end{tikzpicture}

    \vspace{4pt}
    {\small\textbf{Figura Anexo 2:} Diagrama de Comportamiento: Tiempo de respuesta del sistema SIGE (Control de calidad del desempeño).\par}
\end{figure}

\vspace{0.1cm}
\begin{table}[H]
\centering
\footnotesize
\renewcommand{\arraystretch}{1.1}
\begin{tabularx}{\textwidth}{|c|c|X||c|c|X|}
\hline
\rowcolor{tableheader}
\textbf{Sem.} & \textbf{ms} & \textbf{Causa Técnica / Hito} & \textbf{Sem.} & \textbf{ms} & \textbf{Causa Técnica / Hito} \\
\hline
\textbf{S1} & 520 & Despliegue inicial Docker (PostgreSQL/Express) & \textbf{S7} & 470 & Refactorización y sanitización de esquemas \\
\hline
\textbf{S2} & 480 & Esquema relacional y endpoints de usuarios & \textbf{S8} & 350 & Caché en memoria de roles y permisos \\
\hline
\textbf{S3} & 610 & Carga inicial de datos y consultas no indexadas & \textbf{S9} & 300 & Compresión HTTP gzip y optimización JSON \\
\hline
\textbf{S4} & 430 & Indexación de llaves foráneas en tablas SGC & \textbf{S10} & 420 & Pruebas de carga con 20 clientes concurrentes \\
\hline
\textbf{S5} & 390 & Pool de conexiones persistentes PostgreSQL & \textbf{S11} & 380 & Ajuste de índices B-Tree en vistas docentes \\
\hline
\textbf{S6} & 560 & Middlewares de autenticación JWT y RBAC & \textbf{S12} & 360 & Rendimiento estable para ciclo de pruebas QA \\
\hline
\end{tabularx}
\vspace{4pt}
{\small\textbf{Tabla Anexo 2:} Registro de latencia semanal y trazabilidad técnica del Diagrama de Comportamiento.\par}
\end{table}

\clearpage
\subsection*{Anexo 3: Gráfico de Trabajo Completado (Burnup Chart)}
\noindent
El Gráfico de Trabajo Completado (Burnup Chart) refleja el avance global del proyecto SIGE respecto a la Línea Base oficial de 36 tareas de la EDT. Se identifica el estado a la fecha de corte (12 de octubre de 2026) con 22 tareas culminadas (61.1\%) y la proyección hacia la meta final programada para el 30 de octubre de 2026:

\vspace{0.2cm}
\begin{figure}[H]
    \centering
    \begin{tikzpicture}
        \begin{axis}[
            width=0.98\textwidth,
            height=7.2cm,
            xlabel={\textbf{Fases del Proyecto SIGE (Línea Base del Alcance y Cronograma)}},
            ylabel={\textbf{Tareas Acumuladas (EDT)}},
            symbolic x coords={F1,F2,F3,F4,F5},
            xtick=data,
            xticklabels={%
                {\textbf{F1: Análisis}\\[2pt]\scriptsize 12 -- 24 Ago},%
                {\textbf{F2: Diseño}\\[2pt]\scriptsize 25 Ago -- 04 Sep},%
                {\textbf{F3: Core [Hoy]}\\[2pt]\scriptsize 07 Sep -- 08 Oct},%
                {\textbf{F4: QA [Prog.]}\\[2pt]\scriptsize 12 -- 23 Oct},%
                {\textbf{F5: Cierre [Prog.]}\\[2pt]\scriptsize 26 -- 30 Oct}%
            },
            x tick label style={font=\footnotesize, align=center},
            ymin=0, ymax=43,
            ytick={0,6,12,18,24,30,36},
            ymajorgrids=true,
            grid style={dashed, gray!30},
            legend pos=north west,
            legend style={font=\scriptsize, fill=white, fill opacity=0.95, draw=gray!40, rounded corners=2pt}
        ]
            % Zona sombreada de trabajo completado (F1 a F3)
            \fill[green!12, opacity=0.7] (axis cs:F1,0) -- (axis cs:F1,6) -- (axis cs:F2,12) -- (axis cs:F3,22) -- (axis cs:F3,0) -- cycle;
            
            % Línea vertical del corte actual (12 de octubre)
            \draw[thick, dashed, color=gray!70] (axis cs:F3,0) -- (axis cs:F3,38);
            \node[fill=white, draw=gray!60, rounded corners=2pt, font=\tiny\bfseries, text=gray!80, inner sep=1.5pt, anchor=south] at (axis cs:F3, 38) {CORTE: 12 OCT};

            % 1. Línea Base del Alcance (36 Tareas Fijas)
            \addplot[very thick, dashed, color=red!80!black] coordinates {(F1,36) (F2,36) (F3,36) (F4,36) (F5,36)};

            % 2. Trabajo Planificado Acumulado
            \addplot[thick, color=blue!70!black, mark=square*, mark size=3pt, mark options={fill=blue!30, draw=blue!80!black}] coordinates {(F1,6) (F2,12) (F3,25) (F4,31) (F5,36)};
            \node[font=\scriptsize\bfseries, text=blue!80!black, above=3pt] at (axis cs:F3, 25) {25 (Plan)};
            \node[font=\scriptsize\bfseries, text=blue!80!black, above=3pt] at (axis cs:F4, 31) {31 (Plan)};

            % 3. Trabajo Real Completado (Aprobado formalmente)
            \addplot[line width=2.2pt, color=green!60!black, mark=*, mark size=4pt, mark options={fill=green!30, draw=green!70!black}] coordinates {(F1,6) (F2,12) (F3,22)};
            \node[font=\scriptsize\bfseries, text=green!50!black, below=4pt] at (axis cs:F1, 6) {6};
            \node[font=\scriptsize\bfseries, text=green!50!black, below=4pt] at (axis cs:F2, 12) {12};
            \node[fill=green!15, draw=green!60!black, rounded corners=2pt, font=\scriptsize\bfseries, text=green!60!black, inner sep=2pt, below=5pt] at (axis cs:F3, 22) {\textbf{22 compl. (61.1\%)}};

            % 4. Proyección hacia el Cierre
            \addplot[thick, dashed, color=green!60!black, mark=o, mark size=3.5pt, mark options={fill=white, draw=green!70!black}] coordinates {(F3,22) (F4,28) (F5,36)};
            \node[font=\scriptsize\bfseries, text=green!60!black, below=4pt] at (axis cs:F4, 28) {28 (est.)};
            \node[font=\scriptsize\bfseries, text=green!60!black, above=4pt] at (axis cs:F5, 36) {36 (Meta Final)};

            \legend{Zona Trabajo Completado, Línea Base Alcance (36 Tareas), Trabajo Planificado Acumulado, Trabajo Real Concluido (22 Tareas), Proyección Estimada (Fases 4 y 5)}
        \end{axis}
    \end{tikzpicture}
    \vspace{4pt}
    {\small\textbf{Figura Anexo 3:} Gráfico de Trabajo Completado (Burnup Chart) de las 5 Fases de la Línea Base EDT.\par}
\end{figure}

\vspace{0.2cm}
\begin{table}[H]
\centering
\footnotesize
\renewcommand{\arraystretch}{1.15}
\begin{tabularx}{\textwidth}{|c|>{\raggedright\arraybackslash}X|c|c|c|>{\raggedright\arraybackslash}m{3.2cm}|}
\hline
\rowcolor{tableheader}
\textbf{Fase} & \textbf{Denominación Oficial EDT} & \textbf{Período} & \textbf{Tareas} & \textbf{Avance Acum.} & \textbf{Estado Oficial} \\
\hline
\textbf{F1} & Análisis y Planificación Inicial & 12 ago - 24 ago & 6 & \textbf{6 (16.7\%)} & Cumplido (Aprobado) \\
\hline
\textbf{F2} & Diseño Rápido y Arquitectura C4 & 25 ago - 04 sep & 6 & \textbf{12 (33.3\%)} & Cumplido (Aprobado) \\
\hline
\textbf{F3} & Desarrollo Iterativo Core \& IA & 07 sep - 08 oct & 13 & \textbf{22 (61.1\%)} & En ejecución (88.5\% fase) \\
\hline
\textbf{F4} & Pruebas QA, Seguridad y Release & 12 oct - 23 oct & 6 & \textit{28 (77.8\%)} & Programado (Octubre) \\
\hline
\textbf{F5} & Despliegue en VPS y Cierre Formal & 26 oct - 30 oct & 5 & \textit{36 (100\%)} & Programado (Octubre) \\
\hline
\multicolumn{3}{|l|}{\textbf{TOTAL CONSOLIDADO PROYECTO SIGE}} & \textbf{36} & \textbf{22 compl.} & \textbf{En cronograma} \\
\hline
\end{tabularx}
\vspace{4pt}
{\small\textbf{Tabla Anexo 3:} Trazabilidad cuantitativa por fases de la Línea Base del Alcance.\par}
\end{table}

\end{document}